\section{Background: how exact attention is made affordable}\label{sec:background}

This section places the project for a reader who has not met the two techniques it
builds on. Both are \emph{exact}: they reorganise the computation of
\cref{eq:attention} (\cref{sec:algorithm}) without changing its result.

\paragraph{Two walls: time and memory.} For $N$ tokens one attention head forms
$N^2$ query--key scores. The straightforward implementation stores them: at
$N=131{,}072$ the score matrix of a single head has $1.7\times10^{10}$ entries, 64\,GiB
in single precision, more than most accelerators hold. Context windows of language
models have grown from a few thousand tokens to hundreds of thousands, so the
$N\times N$ matrix became a wall long before the $N^2$ arithmetic did.

\paragraph{Inside one device: FlashAttention.} FlashAttention~\cite{flash} removes
the memory wall on a single GPU. It never materialises the score matrix. Queries
and keys are processed in \emph{tiles} small enough for the fast on-chip memory;
the scores of a tile are computed, used and discarded; and the tiles of one query
are combined with the online softmax of \cref{sec:algorithm}, which needs only a
running maximum, a running sum and a running output per query. Memory becomes
linear in $N$, and the kernel is faster because it moves far less data between the
GPU's main memory and its on-chip memory. Later versions improved the division of
work on the GPU~\cite{flash2} and exploited newer hardware~\cite{flash3}. The
approach is now the default way to compute attention: for example, the standard
attention function of PyTorch selects a FlashAttention implementation when one is
available~\cite{sdpa}. In this project the tuned CPU kernel (\cref{sec:serial}) and
the GPU kernel K2 (\cref{sec:gpu}) are hand-written kernels of this kind, tiled and
based on the online softmax; K1, which stores the scores, is the contrast. No
FlashAttention code or other library is used.

\paragraph{Across devices: sequence parallelism.} Tiling fixes the memory of the
attention kernel, but the queries, keys, values and outputs of a whole sequence
must still fit on one device. When they do not, the sequence itself is divided
among devices. Two designs are in use.
\begin{itemize}
\item \textbf{Ring Attention}~\cite{ring} gives each device one block of the
sequence. A device keeps its queries and lets the key/value blocks of all devices
pass by, one per step, around a ring; the online softmax absorbs them in whatever
order they arrive. The context length that can be handled grows with the number of
devices. Two refinements matter for this report. With a causal mask the devices
holding late tokens have far more work than those holding early ones; Striped
Attention~\cite{striped} removes the imbalance by assigning tokens to devices in
stripes instead of contiguous blocks. NVIDIA's training framework implements the
ring as ``context parallelism'' with FlashAttention kernels on each GPU, and
balances the causal work by giving each device one early and one late chunk of the
sequence~\cite{megatroncp}; that assignment is the zigzag partition used in
\cref{sec:causal}.
\item \textbf{DeepSpeed-Ulysses}~\cite{ulysses} also splits the sequence, but before
each attention layer it redistributes the data with an all-to-all exchange so that
every device holds the \emph{whole} sequence for a subset of the heads, computes
those heads locally, and redistributes back. The volume each device sends stays
constant when sequence length and device count grow together, but the number of
devices cannot exceed the number of heads, and each device must hold a full-length
sequence for its heads. The two
designs can be combined on a two-dimensional grid of devices~\cite{usp}.
\end{itemize}

\paragraph{What this project takes from them, and what it adds.} I implement the
ring family: sequence sharding, the online softmax, the ring exchange, a
one-sided alternative, the balanced causal partition, and tiled kernels for CPU and
GPU. The setting differs from the papers in one respect that shapes the results.
There a rank is a GPU that owns tens of thousands of tokens; here a rank is usually
a CPU core that owns a few hundred to a few thousand, which is exactly the regime
where communication starts to matter and can be measured. The published work
reports end-to-end training throughput; this report isolates the attention operator
and asks the parallel-computing questions underneath it: what one core can do,
what a full node costs, when communication shows, whether ``overlapped'' and
``one-sided'' transfers really overlap, and how the answers change when the local
kernel becomes a GPU. Ulysses is not implemented; its all-to-all pattern and its
limit of one device per head (the scaling runs here use a single head) make it a
different experiment.
